Celem niniejszego rozdziału jest zbadanie tworzenia tożsamych programów obiektowych o różnych współczynnikach sprzężenia.
Takie narzędzie jest konieczne do kontroli tej zmiennej niezależnej eksperymentu.

W tym rozdziale zostanie przedstawiona definicja tożsamych programów w rozumieniu klasycznej semantyki operacyjnej,
a następnie jej dostosowanie do kontekstu programowania obiektowego.
Przybliżony zostanie sposób redukcji siły sprzężenia pomiędzy klasami za pomocą powszechnie uznanych metod
refaktoryzacji kodu źródłowego.
Następnie w drodze badania możliwości manipulacji zmienną niezależną zostanie przedstawiona metoda celowego zwiększenia
przypadkowej złożoności: odwróconej refaktoryzacji.
Przedstawione zostaną przykłady takiego zabiegu porównane za pomocą miar sprzężenia wspomnianych w rozdziale~\ref{cha:coupling}.
Na zakończenie rozdziału zostanie omówiony przykład programu \texttt{JavaBus}, użytego do eksperymentu, którego
metodę badawczą przedstawiono w rozdziale~\ref{cha:metoda}, a wyniki w rozdziale~\ref{cha:wyniki}.
